\chapter{Thiết kế vật lý}

\section{Thiết kế lưu trữ vật lý}

Thiết kế lưu trữ vật lý là giai đoạn chuyển đổi mô hình quan hệ logic (gồm 21 quan hệ đã chuẩn hóa đạt BCNF ở Chương 3) thành cấu trúc lưu trữ cụ thể trên hệ quản trị cơ sở dữ liệu Microsoft SQL Server 2022. Nội dung phần này phân tích cách các bản ghi được lưu trữ vật lý dưới dạng các dòng (Rows) trong các trang dữ liệu (Data Pages $8\text{ KB}$), ước lượng không gian lưu trữ và hệ số chặn dựa trên quy mô dữ liệu mẫu thực tế của dự án (\textit{Platform Connecting the Film Photography Community with Darkroom and Studio Services}), đồng thời đánh giá và lựa chọn phương thức tổ chức tập tin phù hợp với đặc thù workload của từng nhóm nghiệp vụ.

\subsection{Khối và bản ghi}

\subsubsection{Cơ chế ánh xạ và đơn vị lưu trữ trong SQL Server}

Trong môi trường Microsoft SQL Server, các khái niệm ở mức logic được hiện thực hóa thành các cấu trúc vật lý cụ thể:
\begin{itemize}
    \item Mỗi \textbf{quan hệ} (Relation) được tạo thành một \textbf{bảng} (Table) trong cơ sở dữ liệu \texttt{FilmPhotographyDB}.
    \item Mỗi \textbf{bản ghi} (Record/Tuple) tương ứng với một \textbf{dòng dữ liệu} (Data Row).
    \item \textbf{Khối dữ liệu (Data Page):} SQL Server sử dụng trang dữ liệu có kích thước cố định $B = 8\text{ KB}$ ($8{,}192\text{ bytes}$) làm đơn vị lưu trữ và cấp phát cơ bản. Trong đó:
    \begin{itemize}
        \item $96\text{ bytes}$ đầu trang dành cho phần tiêu đề (\textit{Page Header}), lưu trữ các thông tin điều khiển như mã định danh trang (\texttt{PageID}), loại trang, dung lượng trống và con trỏ liên kết hai chiều.
        \item Cuối trang là mảng con trỏ dòng (\textit{Row Offset Array}), mỗi dòng chiếm $2\text{ bytes}$ để chỉ định vị trí bắt đầu của dòng tương ứng trong trang.
        \item Không gian khả dụng tối đa cho dữ liệu dòng thực tế trên mỗi trang là $B_{\text{usable}} \approx 8{,}060\text{ bytes}$.
    \end{itemize}
    \item \textbf{Extent ($64\text{ KB}$):} Tám trang dữ liệu liên tiếp ($8 \times 8\text{ KB} = 64\text{ KB}$) tạo thành một Extent. SQL Server quản lý và cấp phát không gian theo Extent để tối ưu hóa hiệu năng đọc/ghi tuần tự trên ổ đĩa.
    \item Khi thực thi bất kỳ truy vấn nào, hệ quản trị cơ sở dữ liệu luôn nạp nguyên trang dữ liệu $8\text{ KB}$ từ đĩa vào vùng nhớ đệm (\textit{Buffer Pool}) thay vì truy xuất từng dòng riêng lẻ.
\end{itemize}

\subsubsection{Phân nhóm các bảng của dự án theo đặc điểm kích thước bản ghi}

Dựa trên cấu trúc 21 bảng trong lược đồ quan hệ đã xác định ở Chương 3, các bảng của dự án được phân thành 4 nhóm kích thước và hành vi lưu trữ:

\begin{enumerate}
    \item \textbf{Nhóm bảng có kích thước bản ghi nhỏ gọn (Compact Rows):}
    Gồm \texttt{USER}, \texttt{PAYMENT}, \texttt{SERVICE\_SESSION}, \texttt{REVIEW}, và \texttt{WORKSHOP\_REGISTRATION}. Các bảng này chủ yếu chứa các thuộc tính số nguyên \texttt{INT} ($4\text{ bytes}$), dữ liệu thời gian \texttt{DATETIME2(0)} ($6\text{ bytes}$), số tiền \texttt{DECIMAL} ($9\text{ bytes}$) và các chuỗi ngắn mã hóa. Kích thước bản ghi trung bình dao động từ $65\text{ bytes}$ đến $190\text{ bytes}$, cho phép mỗi trang $8\text{ KB}$ chứa từ 40 đến hơn 120 dòng.

    \item \textbf{Nhóm bảng có bản ghi lớn do chứa nhiều văn bản mô tả (Wide/Descriptive Rows):}
    Gồm \texttt{CREATIVE\_SPACE}, \texttt{RESOURCE}, \texttt{COMMUNITY\_CONTENT}, \texttt{WORKSHOP}, và \texttt{COMPLAINT}. Do tính chất nền tảng nhiếp ảnh analog cần lưu trữ thông số kỹ thuật (điều kiện ánh sáng, hệ thống thông gió, đặc tính âm học, hóa chất tương thích, nội dung bài viết kỹ thuật), các cột dạng \texttt{NVARCHAR} và \texttt{NVARCHAR(MAX)} làm kích thước bản ghi trung bình tăng lên $300\text{ bytes} - 550\text{ bytes}$, khiến mỗi trang chỉ chứa được khoảng $14 - 25$ dòng.

    \item \textbf{Nhóm bảng giao dịch phát sinh thường xuyên (High-Volume Transaction Rows):}
    Gồm \texttt{RESERVATION}, \texttt{PAYMENT}, \texttt{SERVICE\_SESSION}, \texttt{REVIEW}, và \texttt{COMPLAINT}. Trong đó, \texttt{RESERVATION} là bảng trung tâm phát sinh bản ghi liên tục khi Photographer đặt phòng tối và dịch vụ. Dòng dữ liệu được thiết kế tối giản để giảm số lượng trang vật lý cần đọc khi quét theo khoảng thời gian.

    \item \textbf{Nhóm bảng liên kết quan hệ nhiều--nhiều (Compact Association Rows):}
    Gồm 6 bảng trung gian: \texttt{PACKAGE\_SPACE}, \texttt{PACKAGE\_RESOURCE}, \texttt{RESERVATION\_SPACE}, \texttt{RESERVATION\_RESOURCE}, \texttt{SESSION\_RESOURCE}, và \texttt{WORKSHOP\_REGISTRATION}. Mỗi bản ghi chỉ gồm các khóa ngoại \texttt{INT} và thuộc tính số lượng \texttt{quantity} (\texttt{INT}), có kích thước rất nhỏ ($18\text{ bytes} - 24\text{ bytes}$), đạt hệ số chặn từ $300 - 450\text{ dòng/trang}$, tối ưu cho các phép kết nối bảng (\textit{JOIN}).
\end{enumerate}

\subsection{Ước lượng số khối dữ liệu}

\subsubsection{Mô hình ước lượng lý thuyết}

Trong lý thuyết cơ sở dữ liệu, việc ước lượng số lượng khối/trang dữ liệu cần thiết để lưu trữ một quan hệ $i$ được tính qua các công thức:

\begin{itemize}
    \item \textbf{Hệ số chặn (Blocking Factor - $bfr_i$):} Số lượng bản ghi tối đa có thể lưu trữ trong một trang $8\text{ KB}$ ($B_{\text{usable}} = 8{,}060\text{ bytes}$):
    \[
    bfr_i = \left\lfloor \frac{B_{\text{usable}}}{R_i} \right\rfloor
    \]
    với $R_i$ là kích thước trung bình của một dòng trong quan hệ $i$ (bao gồm chi phí quản lý dòng \textit{Row Header} $4\text{ bytes}$, \textit{Null Bitmap} và \textit{Column Offset Array} $\approx 14 - 22\text{ bytes}$).

    \item \textbf{Số lượng trang dữ liệu lý thuyết ($b_i$):} Với quan hệ có $r_i$ bản ghi:
    \[
    b_i = \left\lceil \frac{r_i}{bfr_i} \right\rceil
    \]

    \item \textbf{Tổng dung lượng lưu trữ ước lượng ($S_i$):}
    \[
    S_i = b_i \times 8\text{ KB}
    \]
\end{itemize}

\subsubsection{Bảng ước lượng lưu trữ chi tiết cho 21 quan hệ của dự án}

Dựa trên số lượng bản ghi thực tế của bộ dữ liệu mẫu đã được nạp trong cơ sở dữ liệu SQL Server của dự án ($5{,}710$ bản ghi), bảng thống kê ước lượng lưu trữ vật lý được trình bày dưới đây:

\begin{table}[htbp]
\centering
\caption{Ước lượng lưu trữ dữ liệu vật lý cho 21 bảng trong SQL Server}
\label{tab:storage_estimation}
\renewcommand{\arraystretch}{1.2}
\scriptsize
\setlength{\tabcolsep}{3pt}
\begin{tabular}{|c|l|c|r|r|r|r|r|}
\hline
\textbf{STT} & \textbf{Tên bảng / Quan hệ} & \textbf{Nhóm dữ liệu} & \textbf{Số bản ghi ($r$)} & \textbf{$R$ trung bình} & \textbf{$bfr$} & \textbf{Số page ($b$)} & \textbf{Dung lượng (KB)} \\
\hline
1 & \texttt{USER} & Người dùng & 100 & 190 bytes & 42 & 3 & 24 KB \\
2 & \texttt{SERVICE\_PROVIDER} & Nhà cung cấp & 20 & 330 bytes & 24 & 1 & 8 KB \\
3 & \texttt{CREATIVE\_SPACE} & Không gian sáng tạo & 60 & 550 bytes & 14 & 5 & 40 KB \\
4 & \texttt{RESOURCE} & Thiết bị, vật tư & 150 & 320 bytes & 25 & 6 & 48 KB \\
5 & \texttt{MAINTENANCE} & Lịch sử bảo trì & 100 & 200 bytes & 40 & 3 & 24 KB \\
6 & \texttt{SERVICE\_PACKAGE} & Gói dịch vụ & 60 & 240 bytes & 33 & 2 & 16 KB \\
7 & \texttt{PROMOTION} & Khuyến mãi & 40 & 230 bytes & 35 & 2 & 16 KB \\
8 & \texttt{RESERVATION} & Giao dịch đặt chỗ & 500 & 80 bytes & 100 & 5 & 40 KB \\
9 & \texttt{PAYMENT} & Giao dịch thanh toán & 450 & 110 bytes & 73 & 7 & 56 KB \\
10 & \texttt{SERVICE\_SESSION} & Phiên sử dụng & 400 & 65 bytes & 124 & 4 & 32 KB \\
11 & \texttt{REVIEW} & Đánh giá dịch vụ & 300 & 170 bytes & 47 & 7 & 56 KB \\
12 & \texttt{COMMUNITY\_CONTENT} & Bài viết cộng đồng & 300 & 390 bytes & 20 & 15 & 120 KB \\
13 & \texttt{WORKSHOP} & Buổi workshop & 50 & 380 bytes & 21 & 3 & 24 KB \\
14 & \texttt{PHOTO} & Bộ sưu tập ảnh & 500 & 270 bytes & 29 & 18 & 144 KB \\
15 & \texttt{COMPLAINT} & Khiếu nại dịch vụ & 50 & 380 bytes & 21 & 3 & 24 KB \\
16 & \texttt{PACKAGE\_SPACE} & Gói -- Không gian (M:N) & 120 & 18 bytes & 447 & 1 & 8 KB \\
17 & \texttt{PACKAGE\_RESOURCE} & Gói -- Tài nguyên (M:N) & 180 & 22 bytes & 366 & 1 & 8 KB \\
18 & \texttt{RESERVATION\_SPACE} & Đặt chỗ -- Không gian (M:N) & 600 & 18 bytes & 447 & 2 & 16 KB \\
19 & \texttt{RESERVATION\_RESOURCE} & Đặt chỗ -- Tài nguyên (M:N) & 800 & 22 bytes & 366 & 3 & 24 KB \\
20 & \texttt{SESSION\_RESOURCE} & Phiên -- Tài nguyên (M:N) & 600 & 22 bytes & 366 & 2 & 16 KB \\
21 & \texttt{WORKSHOP\_REGISTRATION} & Đăng ký Workshop (M:N) & 300 & 24 bytes & 335 & 1 & 8 KB \\
\hline
\multicolumn{3}{|c|}{\textbf{Tổng cộng toàn hệ thống}} & \textbf{5,680} & -- & -- & \textbf{94 pages} & \textbf{752 KB} \\
\hline
\end{tabular}
\end{table}

\subsubsection{Phân biệt ước lượng lý thuyết và lưu trữ thực tế trong SQL Server}

Cần phân biệt rõ giữa mô hình tính toán lý thuyết và cơ chế lưu trữ thực tế của SQL Server Engine:
\begin{itemize}
    \item \textbf{Cơ chế cấp phát Extent ($64\text{ KB}$):} SQL Server cấp phát không gian theo từng Extent (8 trang liên tiếp). Do đó, dung lượng cấp phát vật lý thực tế trên đĩa (\textit{Reserved Space}) thường lớn hơn dung lượng dữ liệu hữu ích (\textit{Data Space}) do chứa các trang đệm chưa dùng hết.
    \item \textbf{Cột có độ dài biến đổi (\texttt{NVARCHAR}, \texttt{NVARCHAR(MAX)}):} Khi kích thước văn bản vượt quá $8\text{ KB}$, SQL Server tự động đẩy phần vượt quá sang các trang riêng biệt (\textit{Row-Overflow Pages} hoặc \textit{LOB Data Pages}) và lưu con trỏ tham chiếu $24\text{ bytes}$ tại trang dữ liệu gốc.
    \item \textbf{Không gian lưu trữ chỉ mục (Index Space):} Ngoài các trang dữ liệu thuần túy, hệ thống còn cấp phát các trang trung gian để xây dựng cấu trúc cây B-Tree cho các chỉ mục Clustered và Non-Clustered.
\end{itemize}

Để đo lường dung lượng và số trang thực tế của từng bảng trong SQL Server, câu lệnh T-SQL sau được thực thi:

\begin{lstlisting}[language=SQL, caption={Truy vấn đo lường lưu trữ vật lý thực tế trong SQL Server}]
USE FilmPhotographyDB;
GO

SELECT 
    t.name AS [TableName],
    p.rows AS [RowCount],
    SUM(a.total_pages) AS [TotalPages],
    SUM(a.used_pages) AS [UsedPages],
    SUM(a.data_pages) AS [DataPages],
    (SUM(a.total_pages) * 8) AS [ReservedSpaceKB],
    (SUM(a.data_pages) * 8) AS [DataSpaceKB]
FROM sys.tables t
INNER JOIN sys.indexes i ON t.object_id = i.object_id
INNER JOIN sys.partitions p ON i.object_id = p.object_id AND i.index_id = p.index_id
INNER JOIN sys.allocation_units a ON p.partition_id = a.container_id
WHERE t.is_ms_shipped = 0
GROUP BY t.name, p.rows
ORDER BY [DataSpaceKB] DESC;
GO
\end{lstlisting}

\subsection{Lựa chọn tổ chức tập tin}

\subsubsection{Phân tích đặc điểm Workload của dự án}

Việc lựa chọn tổ chức tập tin cho từng bảng dựa trên tần suất và đặc thù các luồng truy xuất dữ liệu trong nền tảng:

\begin{enumerate}
    \item \textbf{Tra cứu thông tin người dùng và nhà cung cấp (\texttt{USER}, \texttt{SERVICE\_PROVIDER}):}
    Chủ yếu là các truy vấn tìm kiếm chính xác theo điều kiện bằng trên khóa chính (\texttt{user\_id}, \texttt{provider\_id}) hoặc khóa ứng viên (\texttt{email}, \texttt{phone}). Tần suất đọc chiếm ưu thế so với thao tác ghi.

    \item \textbf{Tra cứu dịch vụ của nhà cung cấp (\texttt{SERVICE\_PROVIDER}, \texttt{CREATIVE\_SPACE}, \texttt{RESOURCE}, \texttt{SERVICE\_PACKAGE}, \texttt{PROMOTION}):}
    Đặc trưng bởi thao tác lọc đa thuộc tính (\textit{Multi-attribute filtering}): tìm phòng tối theo loại (\texttt{space\_type}), sức chứa (\texttt{capacity}), khoảng giá (\texttt{pricing}) và trạng thái; lọc thiết bị theo nhà cung cấp và loại tài nguyên (\texttt{resource\_type}).

    \item \textbf{Giao dịch đặt và sử dụng dịch vụ (\texttt{RESERVATION}, \texttt{PAYMENT}, \texttt{SERVICE\_SESSION}):}
    Nhóm bảng có khối lượng thao tác ghi (Insert/Update) cao nhất. Đặc biệt, \texttt{RESERVATION} liên tục thực hiện truy vấn kiểm tra trùng lịch theo khoảng thời gian (\textit{Time-range queries}: $\texttt{start\_time} \le T \le \texttt{end\_time}$) để tránh xung đột tài nguyên. Bảng \texttt{PAYMENT} và \texttt{SERVICE\_SESSION} có quan hệ $1:1$ với \texttt{RESERVATION}, đòi hỏi truy xuất tức thời theo \texttt{reservation\_id}.

    \item \textbf{Quản lý bảo trì (\texttt{MAINTENANCE}):}
    Thường xuyên truy vấn theo thời gian dự kiến (\texttt{scheduled\_at}) để kiểm tra danh sách thiết bị đang bảo trì nhằm tạm khóa khả năng cho thuê của tài nguyên.

    \item \textbf{Hoạt động cộng đồng và đánh giá (\texttt{COMMUNITY\_CONTENT}, \texttt{PHOTO}, \texttt{WORKSHOP}, \texttt{WORKSHOP\_REGISTRATION}, \texttt{REVIEW}, \texttt{COMPLAINT}):}
    Đặc trưng bởi thao tác hiển thị danh sách bài viết/ảnh mới nhất sắp xếp theo thời gian (\texttt{created\_at} DESC), tra cứu workshop theo ngày diễn ra (\texttt{start\_time}) hoặc tra cứu các khiếu nại chưa xử lý (\texttt{status = 'Open'}).

    \item \textbf{Các quan hệ nhiều--nhiều (Intermediate Relations):}
    Truy xuất hoàn toàn dựa trên phép kết nối (\textit{Equi-Join}) thông qua các khóa ngoại ghép, ví dụ: kiểm tra danh sách tài nguyên trong gói dịch vụ (\texttt{PACKAGE\_RESOURCE}) hoặc danh sách không gian của một lần đặt chỗ (\texttt{RESERVATION\_SPACE}).
\end{enumerate}

\subsection{Heap, Sorted và Hash}

\subsubsection{So sánh tổng quan ba phương thức tổ chức tập tin kinh điển}

\begin{table}[htbp]
\centering
\caption{So sánh các phương thức tổ chức tập tin và mức độ phù hợp với dự án}
\label{tab:file_org_comparison}
\renewcommand{\arraystretch}{1.3}
\scriptsize
\setlength{\tabcolsep}{4pt}
\begin{tabular}{|p{2.2cm}|p{3.1cm}|p{3.1cm}|p{3.1cm}|}
\hline
\textbf{Tiêu chí} & \textbf{Heap File} & \textbf{Sorted File} & \textbf{Hash File} \\
\hline
\textbf{Thứ tự lưu trữ} &
Không duy trì thứ tự logic cố định; bản ghi mới được chèn vào vị trí trống bất kỳ hoặc cuối tập tin. &
Các bản ghi được sắp xếp vật lý liên tục theo giá trị của một khóa sắp xếp (\textit{Sort Key}). &
Vị trí khối được tính toán trực tiếp qua hàm băm $h(K) \rightarrow \text{Address}$. \\
\hline
\textbf{Tìm kiếm bằng khóa ($= K$)} &
Kém hiệu quả; phải quét tuần tự toàn bộ tập tin ($O(b)$ I/O) nếu không có cấu trúc chỉ mục phụ trợ. &
Rất hiệu quả; có thể áp dụng tìm kiếm nhị phân trên các khối dữ liệu ($O(\log_2 b)$ I/O). &
Tối ưu tuyệt đối; thời gian truy xuất lý thuyết xấp xỉ $O(1)$ I/O nếu phân bố đều. \\
\hline
\textbf{Truy vấn theo khoảng ($K_1 \le X \le K_2$)} &
Không hỗ trợ; bắt buộc phải quét toàn bộ bảng ($O(b)$). &
Cực kỳ tối ưu; chỉ cần tìm điểm bắt đầu rồi đọc tuần tự các khối kế tiếp ($O(\log_2 b + \text{kết quả})$). &
Không hỗ trợ; hàm băm phân tán giá trị ngẫu nhiên, không bảo toàn thứ tự. \\
\hline
\textbf{Chi phí thêm mới (Insert)} &
Rất nhanh ($O(1)$ I/O); chỉ cần tìm trang có sẵn không gian trống (\textit{PFS Page}) và ghi vào. &
Tốn kém; phải dịch chuyển các bản ghi phía sau để chèn đúng vị trí hoặc sử dụng vùng tràn (\textit{Overflow Area}). &
Rất nhanh ($O(1)$ I/O) nếu không xảy ra xung đột địa chỉ băm (\textit{Collision}). \\
\hline
\textbf{Mức phù hợp với dự án} &
Phù hợp cho bảng log ghi chép đơn thuần, bảng có kích thước nhỏ ít bản ghi. &
Cần thiết cho các bảng giao dịch thời gian (\texttt{RESERVATION}, \texttt{WORKSHOP}, \texttt{MAINTENANCE}). &
Chỉ phù hợp về mặt lý thuyết với tra cứu khóa định danh đơn lẻ (\texttt{user\_id}, \texttt{provider\_id}). \\
\hline
\end{tabular}
\end{table}

\subsubsection{Phân tích phương thức Heap trong dự án và cơ chế SQL Server}

\paragraph{a. Đặc điểm lý thuyết:}
Heap File cho phép thao tác ghi dữ liệu đạt tốc độ tối đa vì hệ quản trị không mất chi phí tính toán vị trí hay tái cấu trúc dữ liệu. Tuy nhiên, thao tác đọc dữ liệu có điều kiện bắt buộc phải quét qua toàn bộ các trang dữ liệu (\textit{Table Scan}), gây lãng phí bộ nhớ đệm và băng thông I/O.

\paragraph{b. Hiện thực trong SQL Server:}
Trong SQL Server, một bảng không có Clustered Index được gọi là một \textbf{Heap Table}. SQL Server quản lý Heap thông qua hai cấu trúc trang:
\begin{itemize}
    \item \textbf{Page Free Space (PFS):} Theo dõi tỷ lệ phần trăm không gian còn trống của từng trang để cấp phát nhanh cho các thao tác \texttt{INSERT}.
    \item \textbf{Index Allocation Map (IAM):} Quản lý danh sách tất cả các Extent thuộc về bảng Heap để phục vụ việc quét tuần tự.
\end{itemize}

\paragraph{c. Đánh giá áp dụng cho dự án:}
Mặc dù dữ liệu mẫu hiện tại của dự án có quy mô vừa phải ($5{,}710$ bản ghi, dung lượng ước lượng $\approx 704\text{ KB}$), việc để các bảng dưới dạng Heap thuần túy chỉ chấp nhận được ở các bảng có kích thước cực nhỏ như \texttt{SERVICE\_PROVIDER} ($20\text{ dòng} \approx 1\text{ page}$) hoặc \texttt{PACKAGE\_SPACE} ($120\text{ dòng} \approx 1\text{ page}$). Đối với các bảng nghiệp vụ chính, cấu trúc Heap dễ phát sinh hiện tượng chuyển tiếp bản ghi (\textit{Forwarded Records}) khi các cột \texttt{NVARCHAR} được cập nhật dài hơn, gây phân mảnh dữ liệu.

\subsubsection{Phân tích phương thức Sorted và cơ chế Clustered Index trong SQL Server}

\paragraph{a. Đặc điểm lý thuyết:}
Sorted File mang lại hiệu quả vượt trội cho các truy vấn lọc theo khoảng giá trị hoặc khoảng thời gian. Trong nền tảng nhiếp ảnh phim, các truy vấn cốt lõi đều mang bản chất có thứ tự:
\begin{itemize}
    \item Kiểm tra lịch đặt phòng tối: \texttt{start\_time} và \texttt{end\_time} trong \texttt{RESERVATION}.
    \item Lọc lịch bảo trì thiết bị: \texttt{scheduled\_at} trong \texttt{MAINTENANCE}.
    \item Hiển thị danh sách workshop sắp diễn ra: \texttt{start\_time} trong \texttt{WORKSHOP}.
    \item Dòng thời gian bài viết cộng đồng: \texttt{created\_at} trong \texttt{COMMUNITY\_CONTENT}.
\end{itemize}

\paragraph{b. Hiện thực trong SQL Server:}
Cần nhấn mạnh rằng trong SQL Server disk-based tables, \textbf{không tồn tại một loại tệp riêng biệt gọi là "Sorted File" được cấu hình thủ công}. Thay vào đó, SQL Server hiện thực hóa mô hình dữ liệu có thứ tự thông qua \textbf{Clustered Index (Chỉ mục cụm dựa trên cây B+ Tree)}:
\begin{itemize}
    \item Các trang dữ liệu ở mức lá (\textit{Leaf Level}) của Clustered Index chính là các trang chứa dữ liệu thực tế của bảng.
    \item Các trang lá được liên kết đôi (\textit{Doubly Linked List}) theo đúng thứ tự logic của khóa chỉ mục cụm (\textit{Clustering Key}).
    \item Các mức chỉ mục trung gian (\textit{Root và Intermediate Levels}) tạo thành cây B+ Tree cho phép định vị nhanh trang chứa dữ liệu với chi phí $O(\log_B N)$ I/O.
    \item Khi khai báo khóa chính \texttt{PRIMARY KEY} trong câu lệnh \texttt{CREATE TABLE}, SQL Server tự động tạo Clustered Index trên cột khóa chính này (mặc định), biến bảng thành cấu trúc Clustered Table.
\end{itemize}

\paragraph{c. Đánh giá áp dụng cho dự án:}
Tất cả 21 bảng trong dự án đều được tổ chức dưới dạng Clustered Tables dựa trên khóa chính:
\begin{itemize}
    \item Đối với 15 bảng thực thể chính: Khóa chính dạng \texttt{IDENTITY(1,1)} tăng dần đơn điệu (\texttt{user\_id}, \texttt{reservation\_id},...) giúp loại bỏ hoàn toàn hiện tượng phân tách trang (\textit{Page Split}) khi chèn dữ liệu mới, đảm bảo các trang dữ liệu luôn được lấp đầy tối đa ($\approx 100\%$).
    \item Đối với 6 bảng quan hệ nhiều--nhiều: Khóa chính ghép (ví dụ \texttt{(package\_id, space\_id)}) đóng vai trò khóa cụm, giúp gom cụm các dòng liên quan về mặt vật lý trên cùng một trang dữ liệu, tối ưu hóa triệt để thao tác kết nối bảng.
\end{itemize}

\subsubsection{Phân tích phương thức Hash và đối chiếu với SQL Server}

\paragraph{a. Đặc điểm lý thuyết:}
Hash File sử dụng hàm băm toán học để ánh xạ trực tiếp giá trị khóa thành địa chỉ khối vật lý. Phương thức này cực kỳ tối ưu cho các thao tác tìm kiếm chính xác theo giá trị bằng ($= K$), chẳng hạn như tìm kiếm người dùng theo \texttt{user\_id} hay tìm thiết bị theo \texttt{resource\_id}.

Tuy nhiên, Hash File bộc lộ các nhược điểm chí mạng:
\begin{itemize}
    \item Hoàn toàn không thể hỗ trợ các truy vấn tìm kiếm theo khoảng (Range Queries) hoặc sắp xếp (Order By) --- vốn là yêu cầu sống còn của hệ thống đặt phòng và workshop.
    \item Khi số lượng bản ghi tăng trưởng vượt quá sức chứa ban đầu, hiện tượng xung đột băm gia tăng khiến hệ thống phải duyệt qua các chuỗi khối tràn (\textit{Overflow Buckets}), làm suy giảm hiệu năng nghiêm trọng.
\end{itemize}

\paragraph{b. Hiện thực trong SQL Server:}
Trong hệ quản trị SQL Server truyền thống (Disk-Based Database Engine):
\begin{itemize}
    \item \textbf{Không hỗ trợ tổ chức tập tin dạng Hash Table trên đĩa.} Mọi bảng lưu trữ trên đĩa chỉ có thể là Heap hoặc Clustered Index (B+ Tree).
    \item Cơ chế \textbf{Hash Index} chỉ xuất hiện trong tính năng \textbf{In-Memory OLTP (Memory-Optimized Tables)} của SQL Server, nơi dữ liệu được lưu hoàn toàn trên RAM và sử dụng mảng con trỏ xô (\textit{Hash Bucket Array}) không sử dụng trang $8\text{ KB}$.
\end{itemize}

\paragraph{c. Đánh giá áp dụng cho dự án:}
Do nền tảng sử dụng mô hình cơ sở dữ liệu quan hệ truyền thống trên đĩa (Disk-based SQL Server), cấu trúc Hash File ở mức tổ chức tập tin không được áp dụng trực tiếp. Thay vào đó, nhu cầu tìm kiếm nhanh theo khóa bằng sẽ được giải quyết triệt để thông qua Clustered Index và Non-Clustered B-Tree Indexes.

\subsubsection{Kết luận lựa chọn thiết kế lưu trữ vật lý cho dự án}

Từ các phân tích chuyên sâu ở trên, thiết kế lưu trữ vật lý cho nền tảng \textit{Platform Connecting the Film Photography Community with Darkroom and Studio Services} đi đến các quyết định kỹ thuật thống nhất:

\begin{enumerate}
    \item \textbf{Không áp dụng một phương thức duy nhất cho toàn bộ hệ thống:} Việc tổ chức lưu trữ được phân hóa rõ ràng theo bản chất workload của từng nhóm bảng.
    \item \textbf{Cấu trúc bảng chuẩn hóa trên Clustered B-Tree Table:} Toàn bộ 21 quan hệ được tổ chức dưới dạng bảng có chỉ mục cụm (Clustered Tables) với khóa cụm là khóa chính tăng dần \texttt{IDENTITY} (cho bảng thực thể) hoặc khóa ngoại ghép (cho bảng trung gian), đảm bảo cân bằng tối ưu giữa tốc độ ghi mới và tốc độ truy xuất điểm.
    \item \textbf{Bổ sung Non-Clustered Indexes cho các thuộc tính thời gian và khóa ngoại:} Để bù đắp cho việc khóa cụm đã được đặt trên khóa chính, các truy vấn tìm kiếm theo khoảng thời gian (\texttt{RESERVATION.start\_time}, \texttt{WORKSHOP.start\_time}) và các thao tác kết nối theo nhà cung cấp (\texttt{CREATIVE\_SPACE.provider\_id}, \texttt{RESOURCE.provider\_id}) sẽ được hỗ trợ toàn diện bằng các chỉ mục không cụm (Non-Clustered Indexes) được trình bày chi tiết ở phần tiếp theo.
\end{enumerate}

```latex
\section{Thiết kế chỉ mục}

Thiết kế chỉ mục là một nội dung của thiết kế mức vật lý, được thực hiện
dựa trên lược đồ quan hệ đã được xác định ở thiết kế mức logic. Mục tiêu của
việc thiết kế chỉ mục là giảm chi phí truy xuất dữ liệu đối với các truy vấn
thường xuyên, đặc biệt là các truy vấn tìm kiếm, lọc, sắp xếp và kết nối giữa
các quan hệ.

Trong phạm vi hệ thống, chỉ mục được lựa chọn dựa trên 21 quan hệ đã được
xác định trong lược đồ logic cuối cùng, bao gồm 15 quan hệ chính và 6 quan hệ
trung gian. Các chỉ mục không tạo thêm thực thể hoặc thuộc tính mới mà chỉ
được xây dựng trên các thuộc tính đã tồn tại trong lược đồ logic.

Việc lựa chọn chỉ mục được xem xét dựa trên ba yếu tố chính: cấu trúc khóa
của quan hệ, các thuộc tính thường xuyên được sử dụng trong truy vấn và
workload của hệ thống.

\subsection{Xác định nhu cầu sử dụng chỉ mục}

Nhu cầu sử dụng chỉ mục được xác định từ các nghiệp vụ chính của hệ thống
được mô tả ở Chương 0 và các quan hệ đã chuyển đổi ở Chương 2.

Các nhóm truy xuất chính gồm:

\begin{enumerate}
	\item Truy xuất bản ghi theo khóa chính.
	\item Truy xuất dữ liệu thông qua khóa ngoại.
	\item Tìm kiếm và lọc Creative Space.
	\item Tìm kiếm và lọc Resource.
	\item Tìm kiếm và lọc Service Package.
	\item Truy xuất Reservation theo User và thời gian.
	\item Truy xuất Payment theo Reservation hoặc mã giao dịch.
	\item Truy xuất Service Session theo Reservation.
	\item Truy xuất Community Content theo loại và thời gian.
	\item Truy xuất Workshop theo thời gian.
	\item Truy xuất các bản ghi trong các quan hệ trung gian để thực hiện
	các phép JOIN.
\end{enumerate}

\subsubsection*{a) Truy xuất theo khóa chính}

Mỗi quan hệ chính trong lược đồ logic đều có khóa chính dùng để định danh
duy nhất một bản ghi. Các khóa chính gồm:

\begin{itemize}
	\item \texttt{USER.user\_id}.
	\item \texttt{SERVICE\_PROVIDER.provider\_id}.
	\item \texttt{CREATIVE\_SPACE.space\_id}.
	\item \texttt{RESOURCE.resource\_id}.
	\item \texttt{MAINTENANCE.maintenance\_id}.
	\item \texttt{SERVICE\_PACKAGE.package\_id}.
	\item \texttt{PROMOTION.promotion\_id}.
	\item \texttt{RESERVATION.reservation\_id}.
	\item \texttt{PAYMENT.payment\_id}.
	\item \texttt{SERVICE\_SESSION.session\_id}.
	\item \texttt{REVIEW.review\_id}.
	\item \texttt{COMMUNITY\_CONTENT.content\_id}.
	\item \texttt{WORKSHOP.workshop\_id}.
	\item \texttt{PHOTO.photo\_id}.
	\item \texttt{COMPLAINT.complaint\_id}.
\end{itemize}

Đối với các quan hệ trung gian có khóa chính ghép, chỉ mục trên khóa chính
ghép được sử dụng để hỗ trợ việc xác định nhanh một tổ hợp liên kết:

\begin{itemize}
	\item \texttt{PACKAGE\_SPACE(package\_id, space\_id)}.
	\item \texttt{PACKAGE\_RESOURCE(package\_id, resource\_id)}.
	\item \texttt{RESERVATION\_SPACE(reservation\_id, space\_id)}.
	\item \texttt{RESERVATION\_RESOURCE(reservation\_id, resource\_id)}.
	\item \texttt{SESSION\_RESOURCE(session\_id, resource\_id)}.
	\item \texttt{WORKSHOP\_REGISTRATION(user\_id, workshop\_id)}.
\end{itemize}

\subsubsection*{b) Truy xuất theo khóa ngoại}

Các khóa ngoại thường được sử dụng trong phép JOIN và trong các truy vấn
tìm các bản ghi liên quan đến một đối tượng cụ thể. Do đó, các thuộc tính
khóa ngoại là các ứng viên quan trọng để xây dựng chỉ mục.

Các thuộc tính khóa ngoại chính gồm:

\begin{itemize}
	\item \texttt{CREATIVE\_SPACE.provider\_id}.
	\item \texttt{RESOURCE.provider\_id}.
	\item \texttt{MAINTENANCE.resource\_id}.
	\item \texttt{SERVICE\_PACKAGE.provider\_id}.
	\item \texttt{PROMOTION.provider\_id}.
	\item \texttt{RESERVATION.user\_id}.
	\item \texttt{RESERVATION.package\_id}.
	\item \texttt{PAYMENT.reservation\_id}.
	\item \texttt{SERVICE\_SESSION.reservation\_id}.
	\item \texttt{REVIEW.user\_id}.
	\item \texttt{REVIEW.reservation\_id}.
	\item \texttt{COMMUNITY\_CONTENT.user\_id}.
	\item \texttt{WORKSHOP.organizer\_id}.
	\item \texttt{PHOTO.user\_id}.
	\item \texttt{COMPLAINT.user\_id}.
\end{itemize}

Ngoài ra, các quan hệ trung gian sử dụng các khóa ngoại ghép để liên kết
với các quan hệ chính. Các khóa ngoại này được sử dụng trong các truy vấn
JOIN và truy xuất các đối tượng liên quan.

\subsubsection*{c) Truy xuất theo thuộc tính nghiệp vụ}

Một số thuộc tính không phải khóa được sử dụng trong các thao tác tìm kiếm
và lọc. Các thuộc tính có nhu cầu xem xét chỉ mục gồm:

\begin{itemize}
	\item \texttt{CREATIVE\_SPACE.space\_type}, \texttt{capacity},
	\texttt{status}.
	\item \texttt{RESOURCE.resource\_type}, \texttt{condition},
	\texttt{status}.
	\item \texttt{SERVICE\_PACKAGE.price}, \texttt{status}.
	\item \texttt{RESERVATION.start\_time}, \texttt{end\_time},
	\texttt{status}.
	\item \texttt{PAYMENT.transaction\_code}, \texttt{payment\_status}.
	\item \texttt{COMMUNITY\_CONTENT.content\_type}, \texttt{created\_at}.
	\item \texttt{WORKSHOP.start\_time}, \texttt{status}.
	\item \texttt{COMPLAINT.status}, \texttt{created\_at}.
\end{itemize}

Các thuộc tính trên chỉ là các ứng viên chỉ mục. Việc tạo chỉ mục thực tế
phải dựa trên tần suất truy vấn và kích thước dữ liệu.

\subsection{Chỉ mục Clustered/Nonclustered}

Chỉ mục Clustered và Nonclustered được xem xét dựa trên cách hệ quản trị cơ
sở dữ liệu tổ chức và truy xuất dữ liệu.

Đối với các quan hệ chính, khóa chính có mức độ truy xuất cao và thường được
sử dụng để định danh bản ghi. Vì vậy, khóa chính là ứng viên ưu tiên cho cấu
trúc chỉ mục chính của quan hệ.

Các quan hệ có khóa chính đơn gồm:

\begin{itemize}
	\item \texttt{USER}.
	\item \texttt{SERVICE\_PROVIDER}.
	\item \texttt{CREATIVE\_SPACE}.
	\item \texttt{RESOURCE}.
	\item \texttt{MAINTENANCE}.
	\item \texttt{SERVICE\_PACKAGE}.
	\item \texttt{PROMOTION}.
	\item \texttt{RESERVATION}.
	\item \texttt{PAYMENT}.
	\item \texttt{SERVICE\_SESSION}.
	\item \texttt{REVIEW}.
	\item \texttt{COMMUNITY\_CONTENT}.
	\item \texttt{WORKSHOP}.
	\item \texttt{PHOTO}.
	\item \texttt{COMPLAINT}.
\end{itemize}

Đối với sáu quan hệ trung gian, khóa chính ghép được sử dụng để hỗ trợ
truy xuất theo tổ hợp khóa:

\begin{itemize}
	\item \texttt{PACKAGE\_SPACE}.
	\item \texttt{PACKAGE\_RESOURCE}.
	\item \texttt{RESERVATION\_SPACE}.
	\item \texttt{RESERVATION\_RESOURCE}.
	\item \texttt{SESSION\_RESOURCE}.
	\item \texttt{WORKSHOP\_REGISTRATION}.
\end{itemize}

Chỉ mục Nonclustered phù hợp với các thuộc tính không phải khóa chính nhưng
được sử dụng thường xuyên trong điều kiện tìm kiếm, lọc hoặc JOIN. Các
khóa ngoại như \texttt{provider\_id}, \texttt{user\_id}, \texttt{resource\_id},
\texttt{package\_id} và \texttt{reservation\_id} là những ứng viên quan trọng.

Việc lựa chọn Clustered hay Nonclustered cụ thể phụ thuộc vào hệ quản trị cơ
sở dữ liệu được lựa chọn khi triển khai. Do Task của dự án cho phép sử dụng
PostgreSQL hoặc MySQL, thiết kế chỉ mục ở mức này không cố định cách cài đặt
phụ thuộc vào một DBMS cụ thể.

\subsection{Chỉ mục B+}

B+ tree phù hợp với các truy vấn tìm kiếm theo thứ tự hoặc theo khoảng giá
trị. Trong hệ thống, các thuộc tính thời gian và giá trị số là những ứng
viên phù hợp cho loại chỉ mục này.

\subsubsection*{a) Reservation theo thời gian}

Reservation là quan hệ trung tâm trong nghiệp vụ đặt chỗ. Các thuộc tính
\texttt{start\_time} và \texttt{end\_time} được sử dụng để xác định khoảng
thời gian đặt dịch vụ.

Các truy vấn có thể yêu cầu tìm Reservation trong một khoảng thời gian:

\[
T_1 \leq \texttt{start\_time} < T_2
\]

Do đó, B+ tree trên \texttt{RESERVATION.start\_time} là một ứng viên phù
hợp khi workload có nhiều truy vấn theo khoảng thời gian.

\subsubsection*{b) Service Package theo giá}

Thuộc tính \texttt{SERVICE\_PACKAGE.price} có thể được sử dụng trong các
truy vấn lọc theo khoảng giá:

\[
P_1 \leq \texttt{price} \leq P_2
\]

B+ tree hỗ trợ hiệu quả dạng truy vấn này và có thể được xem xét khi chức
năng tìm kiếm Service Package thường xuyên sử dụng bộ lọc theo giá.

\subsubsection*{c) Workshop theo thời gian}

Thuộc tính \texttt{WORKSHOP.start\_time} có thể được sử dụng để tìm các
Workshop diễn ra trong một khoảng thời gian nhất định hoặc sắp xếp các
Workshop theo thời gian.

Do đó, B+ tree là ứng viên phù hợp cho:

\[
\texttt{WORKSHOP.start\_time}
\]

\subsubsection*{d) Community Content theo thời gian}

Đối với \texttt{COMMUNITY\_CONTENT}, thuộc tính \texttt{created\_at} có thể
được sử dụng để lấy nội dung mới hoặc lọc nội dung theo khoảng thời gian.

Vì vậy, B+ tree trên:

\[
\texttt{COMMUNITY\_CONTENT.created\_at}
\]

có thể được sử dụng khi workload thường xuyên truy xuất nội dung theo thời
gian.

\subsubsection*{e) Đánh giá}

B+ tree được ưu tiên khi workload có:

\begin{itemize}
	\item Truy vấn tìm kiếm theo khoảng.
	\item Truy vấn so sánh lớn hơn hoặc nhỏ hơn.
	\item Truy vấn sắp xếp theo thứ tự.
	\item Truy vấn lấy giá trị nhỏ nhất hoặc lớn nhất.
\end{itemize}

Do đó, B+ tree phù hợp hơn Hash index đối với các thuộc tính
\texttt{start\_time}, \texttt{end\_time}, \texttt{price} và
\texttt{created\_at} khi workload sử dụng các phép so sánh hoặc khoảng giá
trị.

\subsection{Chỉ mục Hash}

Hash index phù hợp với các truy vấn tìm kiếm chính xác theo giá trị, đặc biệt
là các truy vấn có điều kiện:

\[
A = value
\]

Hash index không phù hợp với các truy vấn cần sắp xếp hoặc tìm kiếm theo
khoảng.

Trong hệ thống, một số thuộc tính có thể xem xét Hash index khi workload
chủ yếu là truy xuất chính xác:

\begin{itemize}
	\item \texttt{USER.user\_id}.
	\item \texttt{SERVICE\_PROVIDER.provider\_id}.
	\item \texttt{RESOURCE.resource\_id}.
	\item \texttt{RESERVATION.reservation\_id}.
	\item \texttt{PAYMENT.transaction\_code}.
\end{itemize}

Tuy nhiên, đối với các khóa chính, hệ quản trị cơ sở dữ liệu đã có cơ chế
chỉ mục thông qua ràng buộc khóa chính. Vì vậy, không cần tạo thêm Hash index
trên cùng thuộc tính nếu không có yêu cầu đặc biệt từ workload.

Đối với \texttt{PAYMENT.transaction\_code}, Hash index có thể được xem xét
nếu nghiệp vụ thường xuyên tra cứu một giao dịch bằng mã giao dịch chính xác.

Việc lựa chọn Hash index cần phụ thuộc vào khả năng hỗ trợ của DBMS được
lựa chọn. Nếu workload yêu cầu tìm kiếm theo khoảng hoặc sắp xếp thì B+ tree
phù hợp hơn.

\subsection{Chỉ mục kết hợp}

Chỉ mục kết hợp được xây dựng trên nhiều thuộc tính và được sử dụng khi các
thuộc tính thường xuyên xuất hiện cùng nhau trong điều kiện truy vấn.

Trong hệ thống, các chỉ mục kết hợp được xem xét từ những workload chính.

\subsubsection*{a) Reservation theo User và thời gian}

Một trong những workload quan trọng là truy xuất các Reservation của một
User theo thời gian.

Do đó, có thể xem xét chỉ mục:

\[
(\texttt{RESERVATION.user\_id},
\texttt{RESERVATION.start\_time})
\]

Chỉ mục này phù hợp với các truy vấn bắt đầu bằng điều kiện trên
\texttt{user\_id} và tiếp tục lọc hoặc sắp xếp theo \texttt{start\_time}.

\subsubsection*{b) Creative Space theo trạng thái và loại}

Nếu workload thường xuyên tìm các Creative Space theo trạng thái và loại
không gian, có thể xem xét:

\[
(\texttt{CREATIVE\_SPACE.status},
\texttt{CREATIVE\_SPACE.space\_type})
\]

Chỉ mục này hỗ trợ các truy vấn lọc Creative Space theo hai thuộc tính.

\subsubsection*{c) Resource theo loại và trạng thái}

Đối với Resource, workload tìm kiếm tài nguyên theo loại và trạng thái có
thể được hỗ trợ bởi:

\[
(\texttt{RESOURCE.resource\_type},
\texttt{RESOURCE.status})
\]

Chỉ mục này phù hợp với các truy vấn tìm Resource theo loại tài nguyên và
trạng thái sử dụng.

\subsubsection*{d) Community Content theo loại và thời gian}

Đối với Community Content, nếu workload thường xuyên tìm nội dung theo loại
và thời gian tạo, có thể xem xét:

\[
(\texttt{COMMUNITY\_CONTENT.content\_type},
\texttt{COMMUNITY\_CONTENT.created\_at})
\]

Chỉ mục hỗ trợ việc lọc theo loại nội dung và tiếp tục sắp xếp hoặc lọc theo
thời gian tạo.

\subsubsection*{e) Nguyên tắc lựa chọn thứ tự thuộc tính}

Thứ tự thuộc tính trong chỉ mục kết hợp phải dựa trên workload. Thuộc tính
được sử dụng trong điều kiện lọc chính thường được đặt ở vị trí đầu tiên.

Ví dụ:

\[
(\texttt{user\_id},\texttt{start\_time})
\]

phù hợp với truy vấn có điều kiện:

\[
\texttt{user\_id}=U
\]

sau đó lọc theo \texttt{start\_time}.

Không nên tạo nhiều chỉ mục kết hợp chỉ dựa trên khả năng có thể xuất hiện
trong truy vấn. Mỗi chỉ mục phải có workload tương ứng để tránh tăng chi phí
lưu trữ và chi phí cập nhật dữ liệu.

\subsection{Lựa chọn chỉ mục theo workload}

Dựa trên các nghiệp vụ đã xác định ở Chương 0, cấu trúc quan hệ ở Chương 2
và lược đồ logic cuối cùng ở Chương 3, các workload chính của hệ thống được
tổng hợp như sau:

\begin{enumerate}
	\item Truy xuất trực tiếp theo khóa chính.
	\item JOIN giữa các quan hệ thông qua khóa ngoại.
	\item Tìm kiếm và lọc Creative Space.
	\item Tìm kiếm và lọc Resource.
	\item Tìm kiếm Service Package theo giá và trạng thái.
	\item Truy xuất Reservation của User theo thời gian.
	\item Kiểm tra và truy xuất Reservation theo khoảng thời gian.
	\item Tra cứu Payment theo Reservation hoặc
	\texttt{transaction\_code}.
	\item Truy xuất Service Session theo Reservation.
	\item Truy xuất Community Content theo loại và thời gian.
	\item Truy xuất Workshop theo thời gian.
	\item Truy xuất dữ liệu từ các quan hệ trung gian để phục vụ các phép JOIN.
\end{enumerate}

Bảng sau tổng hợp lựa chọn chỉ mục theo workload:

\begin{table}[H]
	\centering
	\caption{Lựa chọn chỉ mục theo workload}
	\label{tab:index_workload}
	\renewcommand{\arraystretch}{1.2}
	\scriptsize
	\setlength{\tabcolsep}{4pt}
	\begin{tabular}{|p{3.4cm}|p{3.2cm}|p{5.2cm}|}
		\hline
		\textbf{Workload} &
		\textbf{Chỉ mục đề xuất} &
		\textbf{Mục đích} \\
		\hline
		
		Truy xuất theo khóa chính &
		Chỉ mục trên PK &
		Tìm kiếm trực tiếp một bản ghi theo định danh. \\
		\hline
		
		JOIN theo khóa ngoại &
		Nonclustered index trên FK &
		Giảm chi phí truy xuất các bản ghi liên quan khi thực hiện JOIN. \\
		\hline
		
		Reservation theo User và thời gian &
		\texttt{(user\_id,start\_time)} &
		Truy xuất lịch sử Reservation của một User theo thời gian. \\
		\hline
		
		Reservation theo khoảng thời gian &
		B+ tree trên \texttt{start\_time} &
		Hỗ trợ tìm kiếm và lọc Reservation theo thời gian. \\
		\hline
		
		Creative Space theo trạng thái và loại &
		\texttt{(status,space\_type)} &
		Hỗ trợ lọc Creative Space theo trạng thái và loại. \\
		\hline
		
		Resource theo loại và trạng thái &
		\texttt{(resource\_type,status)} &
		Hỗ trợ tìm Resource phù hợp với điều kiện truy vấn. \\
		\hline
		
		Service Package theo giá &
		B+ tree trên \texttt{price} &
		Hỗ trợ tìm kiếm theo khoảng giá. \\
		\hline
		
		Workshop theo thời gian &
		B+ tree trên \texttt{start\_time} &
		Hỗ trợ tìm kiếm và sắp xếp Workshop theo thời gian. \\
		\hline
		
		Community Content theo thời gian &
		B+ tree trên \texttt{created\_at} &
		Hỗ trợ truy xuất nội dung mới hoặc theo khoảng thời gian. \\
		\hline
		
		Payment theo mã giao dịch &
		Hash hoặc B+ tree &
		Hash phù hợp với tìm kiếm chính xác; B+ tree phù hợp nếu cần sắp xếp hoặc
		tìm kiếm theo khoảng. \\
		\hline
		
	\end{tabular}
\end{table}

Việc lựa chọn chỉ mục cuối cùng phải cân bằng giữa hiệu năng đọc và chi phí
ghi. Mỗi chỉ mục bổ sung đều làm tăng dung lượng lưu trữ và có thể làm tăng
chi phí của các thao tác \texttt{INSERT}, \texttt{UPDATE} và
\texttt{DELETE}.

Do đó, các chỉ mục ưu tiên cao của hệ thống là:

\begin{itemize}
	\item Chỉ mục trên khóa chính của các quan hệ.
	\item Chỉ mục trên các khóa ngoại thường xuyên được sử dụng trong JOIN.
	\item Chỉ mục hỗ trợ Reservation theo User và thời gian.
	\item Chỉ mục hỗ trợ truy xuất Reservation theo thời gian.
	\item Chỉ mục phục vụ các truy vấn tìm kiếm Creative Space và Resource
	khi workload đủ lớn.
\end{itemize}

Các chỉ mục trên thuộc tính như \texttt{status}, \texttt{space\_type},
\texttt{resource\_type}, \texttt{content\_type} và các chỉ mục kết hợp chỉ
nên được triển khai khi workload thực tế chứng minh rằng lợi ích truy xuất
lớn hơn chi phí duy trì chỉ mục.

Như vậy, thiết kế chỉ mục được xây dựng hoàn toàn dựa trên lược đồ logic đã
được xác định trước đó. Các quyết định về Clustered/Nonclustered, B+ tree,
Hash và chỉ mục kết hợp là các quyết định ở mức vật lý nhằm tối ưu workload,
không làm thay đổi các quan hệ, thuộc tính, PK hoặc FK của thiết kế logic.
```


\section{Phân vùng dữ liệu}
\subsection{Phân vùng ngang}
\subsection{Phân vùng dọc}
\subsection{Lựa chọn khóa phân vùng}

\section{Tối ưu truy vấn}
\subsection{Phân tích workload}
\subsection{Cây truy vấn và tối ưu phép toán}
\subsection{Ước lượng kích thước trung gian}
\subsection{Ước lượng chi phí truy vấn}
\subsection{Lựa chọn thuật toán kết nối}

\section{Lựa chọn cấu hình vật lý}
\subsection{Bộ đệm (Buffer Pool)}
\subsection{RAID và tổ chức lưu trữ}
\subsection{Row-store và Column-store}
\subsection{Đánh đổi giữa đọc, ghi và dung lượng}

\section{Thiết kế vật lý cuối cùng}
\subsection{Kiểu dữ liệu}
\subsection{Các ràng buộc vật lý}
\subsection{Các chỉ mục}
\subsection{Phân vùng/cấu hình lưu trữ}
\subsection{Tổng hợp thiết kế vật lý}
